\section{Architecture and Primitives}
\label{sec:checkpoints-architecture}

\subsection{Core Objects}
\label{sec:checkpoints-architecture-core-objects}

The checkpoint ledger is built from four core object types. The first is the \emph{verdict document},
a Markdown file named with a checkpoint identifier encoding verdict type and sequence number
(e.g., PROCESS\_INTELLIGENCE\_ALIVE\_001, GGEN\_ECOSYSTEM\_INTEL\_ALIVE\_001). Each verdict
document records the date of sealing, the authority issuing the verdict, the gate criteria evaluated,
the actual counts at assessment time, and the downstream workflows authorized or refused.

The second core object is the \emph{gate criteria record}, encoded in RESEARCH\_CRITERIA.md. This
document defines the gate table before any verdict is sought, making the threshold pre-declared and
falsifiable. The gate table specifies twelve directory-level criteria with minimum file counts and
verification shell commands, plus a commit count threshold (minimum 80 commits, actual 894 at
program ALIVE assessment).

The third core object is the \emph{receipt}, a BLAKE3- or SHA-256-sealed document in the companion
\texttt{receipts/} directory. Receipts are the evidence that verdict documents reference. Each
receipt carries a witness marker (e.g., InductiveWitness, ConformanceWitness, GovernanceWitness)
that binds the receipt to a specific type-law surface.

The fourth core object is the \emph{cryptographic seal}, a hash string appended to each verdict
document. PROCESS\_INTELLIGENCE\_ALIVE\_001 carries a verification status code 0x00 (SUCCESS\_ALIVE).
PI\_RESEARCH\_PROGRAM\_ALIVE\_001 carries a BLAKE3 seal:
\texttt{e7c8f2d94a71b5c3e9f1d6a4b2c8e5f7a1d3c5b7e9f2d4a6c8e0f1a3b5c7d9}.
GGEN\_ECOSYSTEM\_INTEL\_ALIVE\_001 carries an SHA-256 seal:
\texttt{d5d7990b798b31a3962d3bf30f0f3531b26f56a310efc35ad1a89b3f021e85a6}.

\subsection{Admissible Transitions}
\label{sec:checkpoints-architecture-transitions}

The checkpoint ledger enforces a typed state machine over verdict documents. The admissible
transitions are: (1) from an empty ledger to PARTIAL, which requires only that the research foundry
has been bootstrapped and criteria are in progress; (2) from PARTIAL to ALIVE, which requires all
gate criteria to be simultaneously met and all open CRITICAL gaps to have documented remediation
plans; (3) from ALIVE to a specialized ALIVE (e.g., GGEN\_ECOSYSTEM\_INTEL\_ALIVE\_001 sealing after
PROCESS\_INTELLIGENCE\_ALIVE\_001), which requires the parent ALIVE to be sealed and the specialized
gate criteria to pass.

The PARTIAL state is a first-class citizen of this state machine. It is not a transient error to be
hidden; it is a named state with an explicit authorization scope. A PARTIAL verdict authorizes
continuation of research work toward the missing artifacts. It does not authorize downstream
implementation work. This distinction is enforced by the authorization lists in each verdict document.

The decommissioning transition $\delta(P)$ is represented in the Petri net snapshot. The base64-
decoded PETRI\_NET\_SNAPSHOT.md records the final marking: [Decommissioning]:1, [Design]:0,
[Simulation]:0, [Construction]:0, [Actuation]:0. A single token in the terminal place with no
tokens in intermediate places satisfies the proper completion criterion: the process is sound because
it has reached the terminal state from the initial [Design] state with no deadlocks.

\subsection{Boundaries}
\label{sec:checkpoints-architecture-boundaries}

The checkpoint ledger enforces two critical boundaries that are stated explicitly in the ALIVE
verdict documents. The first boundary is the \emph{admissibility boundary}: $R \vdash P_i = \mu(O^*, T, L)$.
No process artifact may be admitted to the manufacturing pipeline without a named witness and a receipt
in the registry. Claims based on unadmitted raw event logs are explicitly not authorized by any
checkpoint.

The second boundary is the \emph{authorization boundary}: downstream workflows are authorized only
when their prerequisites appear in a sealed ALIVE verdict's authorization list. The ALIVE gate
assessment document records five authorized downstream workflows following PROCESS\_INTELLIGENCE\_ALIVE\_001:
board-level claim delivery, M\&A diligence package assembly, wasm4pm graduation initiation, PAPERLAW
downstream workflows, and standards-based conformance claims. It simultaneously records four
unauthorized workflow classes: claims based on unadmitted event logs, conformance scores without a
named witness, M\&A claims without corresponding receipts, and wasm4pm integration without
\texttt{Admission<T, W>} typing at the API boundary.

\subsection{Failure Modes Refused}
\label{sec:checkpoints-architecture-failure-modes}

The checkpoint architecture explicitly refuses four failure modes that the old regime permits. The
first is the \emph{forced ALIVE verdict}: issuing an ALIVE certification before all gate criteria are
met. Gate 12 of the twelve-gate evaluation is titled \textquotedblleft No Forced ALIVE
Verdicts\textquotedblright{} and passes only when the audit confirms zero forced verdicts in the
ledger. The PI\_RESEARCH\_PROGRAM\_ALIVE\_001 document records this gate as passing with the
observation that PROCESS\_INTELLIGENCE\_ALIVE\_001 was sealed only after 12 of 12 gates passed.

The second refused failure mode is the \emph{partial gate treated as passed}. If a gate criterion
requires 8 lifecycle files and the inventory contains 7, the gate fails and the overall verdict is
PARTIAL. There is no partial credit.

The third refused failure mode is the \emph{undocumented gap}. Every open item that cannot be closed
at assessment time must be recorded in the gaps/ directory with a severity classification and a
remediation plan. The ALIVE gate assessment lists five open gaps (GAP\_001 through GAP\_005), each
with a documented timeline and blocking status.

The fourth refused failure mode is the \emph{conformance score without a named witness}. The
ALIVE\_GATE\_ASSESSMENT document records this as an explicit prohibition: conformance scores are not
admissible in board presentations or M\&A diligence packages unless they carry a named witness
marker traceable to a receipt in the registry.
